\section{Exact finite certificate and numerical orientation}
\label{sec:certificate}

The proof of the asymptotic theorem depends on the cited Titchmarsh and
Bombieri--Vinogradov inputs, which no finite computation can certify.  The
algebraic layer, however, admits an exact machine-checkable certificate.  We
also record a separate floating-point diagnostic so that the two roles are
not conflated.

\subsection{Formal certificate}

The program
\texttt{experiments/exact\_zero\_mellin\_certificate.py} represents each
$\log p$ and each product $\log p\log q$ by a formal basis symbol with an
integer coefficient.  It therefore uses no floating approximation to the
von Mangoldt function.  With
\begin{equation}\label{eq:certificate-parameters}
 1000\leq r\leq2000,\qquad h=2,\qquad D=31,
\end{equation}
the committed certificate checks all of the following statements exactly:
\begin{enumerate}[label=\textup{(\alph*)}]
 \item for every one of the $1001$ product values,
 \(
   \sum_{d\mid r}(\Lambda(d)-1)=\log r-\tau(r);
 \)
 \item the direct sum over all $8465$ factor pairs equals the collapsed
 product sum as a formal expression;
 \item the $4035$ pairs with $d\leq31$ and the $4430$ pairs with $d>31$
 recombine to the complete expression;
 \item for every $1\leq d\leq31$, the substitution
 $t=dm+h$ agrees exactly with the corresponding arithmetic-progression
 indexing;
 \item the local Euler factors at $s=0$ used in
 \eqref{eq:phi-Dirichlet-series} agree as rational numbers for every prime
 $p\leq19$.
\end{enumerate}
The resulting complete formal expression has $431$ nonzero basis terms.
The canonical payload SHA--256 is
\begin{center}
 \nolinkurl{6686C9E926CCB05F4DB0BF56A6A8F2811D5DED9A020FB905AC8DC54D350EDAEE},
\end{center}
and the SHA--256 of the indented committed JSON file is
\begin{center}
 \nolinkurl{A0D6303B26118D756406677A748D74160065390E5B4A55F70FDD2F1549E2BE43}.
\end{center}
These hashes identify the finite conventions and artifact.  They do not
certify \cref{thm:ABL-smooth-input,thm:fixed-product-Type-I}, nor are they
used in their proofs.

\subsection{A deliberately non-theorem diagnostic}

For numerical orientation only, the second program replaces the smooth
weight by the sharp box $X<r\leq2X$, takes $h=2$ and
$D=\lfloor\sqrt X\rfloor$, and computes the exact decomposition using
floating logarithms.  The entries in \cref{tab:sharp-box-diagnostic} are
normalized by $X\log X$.  The last column is
$X^{-1}\sum_{X<r\leq2X}\Lambda(r+2)$.

\begin{table}[ht]
\centering
\caption{Sharp-box implementation diagnostic; this is not a verification
of the smooth asymptotic theorem.}
\label{tab:sharp-box-diagnostic}
\begin{tabular}{rrrrrr}
\toprule
$X$ & $D$ & complete & $n\leq D$ & $n>D$ & target mass \\
\midrule
$2^{14}$ & 128  & 0.248166 & 0.039334 & 0.208832 & 0.996607 \\
$2^{16}$ & 256  & 0.264715 & 0.054820 & 0.209896 & 1.001175 \\
$2^{18}$ & 512  & 0.270705 & 0.071123 & 0.199581 & 1.001127 \\
$2^{20}$ & 1024 & 0.277987 & 0.080975 & 0.197012 & 1.000196 \\
\bottomrule
\end{tabular}
\end{table}

For comparison, the predicted leading complete ratio is
$1-C_2=0.352134\ldots$, while the square-root tail leading ratio is
$(1-C_2)/2=0.176067\ldots$.  The modest values of $X$, the discontinuous
weight, lower-order terms, and the unspecified logarithmic loss in the
admissible cutoff prevent this table from testing the theorem.  Its purpose
is narrower: it detects implementation mistakes in the complete/small/tail
bookkeeping and displays the slow finite-scale drift.

Both layers are reproducible using only the Python standard library:
\begin{verbatim}
python experiments/exact_zero_mellin_certificate.py \
  --output experiments/data/exact-zero-mellin-certificate.json
python experiments/test_exact_zero_mellin_certificate.py
python experiments/zero_mellin_diagnostic.py \
  --powers 14,16,18,20 --h 2
python experiments/test_zero_mellin_diagnostic.py
\end{verbatim}

\section{Conclusion and the next analytic interface}
\label{sec:conclusion}

The paper isolates one coordinate in which the full prime fluctuation can
be passed through a prescribed fixed-shift prime target.  Complete
factor-scale aggregation turns $\Lambda-1$ into $\log-\tau$ exactly.  A
proven Titchmarsh theorem evaluates the total, classical
Bombieri--Vinogradov evaluates the short factor range, and subtraction
evaluates the complete long-factor tail.  When its leading coefficient is
nonzero, the $X\log X$ coefficient at a square-root cutoff is one half of
the complete-sum coefficient.

This positive result is deliberately paired with its limitation.  It is a
statement about one scalar scale coordinate, not about the distribution of
the individual layers composing it.  Mellin inversion shows that a dyadic
factor localization requires nonzero scale frequencies, for which the
elementary divisor collapse disappears.  At the endpoint $m=1$, the
missing localized estimate is already the fixed-shift prime-pair problem.

The next defensible layer is therefore not to infer a localized theorem
from the zero mode.  It is to study the nonzero Mellin family in a region
where existing dispersion or Type~II estimates genuinely apply, keeping
the coefficient class and averaging axes explicit.  A successful result
there would quantify how much scale resolution can be recovered.  A
negative result would still identify the exact norm or coefficient
geometry that blocks recovery.  Neither outcome is assumed here.
